%%%%%%%%%%%%%%%%%%%%%%%%%%%%%%%%%
%%\newpage
%\part{The Bishop--Jones theorem}
%%\section{The Bishop--Jones theorem}
%\label{partbishopjones}
%This part will be divided as follows: In Section \ref{sectionmodified}, we motivate and define the \emph{modified} Poincar\'e exponent of a semigroup, which is used in the statement of Theorem \ref{theorembishopjonesmodified}. In Section \ref{sectionbishopjones} we prove Theorem \ref{theorembishopjonesmodified} and deduce Theorem \ref{theorembishopjonesregular} from Theorem \ref{theorembishopjonesmodified}.


%%%%%%%%%%%%%%%%%%%%%%%%%%%%%%%%
\chapter{The modified Poincar\'e exponent} \label{sectionmodified}
%%%%%%%%%%%%%%%%%%%%%%%%%%%%%%%%
In this chapter we define the modified Poincar\'e exponent of a semigroup. We first recall the classical notion of the Poincar\'e exponent, introduced in the Standard Case by A. F. Beardon in \cite{Beardon1}. Although it is usually defined only for groups, the generalization to semigroups is trivial.

%%%%%%%%%%%%%%%%%%%%%%%%%%%%%%%%%%
%\bigskip
\section{The Poincar\'e exponent of a semigroup}
\label{subsectionpoincare}
%%%%%%%%%%%%%%%%%%%%%%%%%%%%%%%%%%

\begin{definition}
\label{definitionpoincareexponent}
Fix $G\prec\Isom(X)$. For each $s\geq 0$, the series
\[
\Sigma_s(G) := \sum_{g\in G}b^{-s\dogo g}
\]
is called the \emph{Poincar\'e series} of the semigroup $G$ in dimension $s$ (or ``evaluated at $s$'') relative to $b$. The number
\[
\delta_G = \delta(G) := \inf\{s\geq 0:\Sigma_s(G) < \infty\}
\]
is called the \emph{Poincar\'e exponent} of the semigroup $G$ relative to $b$. Here, we let $\inf\emptyset = \infty$.
\end{definition}

\begin{remark}
The Poincar\'e series is usually defined with a summand of $e^{-s\dogo g}$ rather than $b^{-s\dogo g}$. The change of exponents here is important because it relates the Poincar\'e exponent to the metric $\Dist = \Dist_{b,\zero}$ defined in Proposition \ref{propositionDist}. In the Standard Case, and more generally for CAT(-1) spaces, we have made the convention that $b = e$ (see \6\ref{standingassumptions2}), so in this case our series reduces to the classical one.
\end{remark}

\begin{remark}
\label{remarkorbitalcounting}
Given $G\prec\Isom(X)$, we may define the \emph{orbital counting function} of $G$ to be the function
\[
\NN_{X,G}(\rho) = \#\{g\in G : \dogo g \leq \rho\}.
\]
The Poincar\'e series may be written as an integral over the orbital counting function as follows:
\begin{equation}
\label{poincareseriesalternate}
\begin{split}
\Sigma_s(G) &= \log(b^s) \sum_{g\in G} \int_{\dogo g}^\infty b^{-s\rho} \;\dee \rho\\
&= \log(b^s) \int_0^\infty b^{-s\rho} \sum_{g\in G}[\dogo g\leq \rho] \;\dee \rho\\
&= \log(b^s) \int_0^\infty b^{-s\rho} \NN_{X,G}(\rho) \;\dee \rho.
\end{split}
\end{equation}
The Poincar\'e exponent is written in terms of the orbital counting function as
\begin{equation}
\label{poincarealternate}
\delta_G = \limsup_{\rho\to \infty} \frac1\rho\log_b \NN_{X,G}(\rho)
\end{equation}
\end{remark}

\begin{definition}
\label{definitiondivergencetype}
A semigroup $G\prec\Isom(X)$ with $\delta_G < \infty$ is said to be of \emph{convergence type} if $\Sigma_{\delta_G}(G) < \infty$. Otherwise, it is said to be of \emph{divergence type}. In the case where $\delta_G = \infty$, we say that the semigroup is neither of convergence type nor of divergence type.
\end{definition}

The most basic question about the Poincar\'e exponent is whether it is finite. For groups, the finiteness of the Poincar\'e exponent is related to strong discreteness:

\begin{observation}\label{observationnotstronglydiscrete}
Fix $G\leq\Isom(X)$. If $G$ is not strongly discrete, then $\delta_G = \infty$.
\end{observation}
\begin{proof}
Fix $\rho > 0$ such that $\#\{g\in G : \dogo g \leq \rho\} = \infty$. Then for all $s\geq 0$ we have
\[
\Sigma_s(G) \geq \sum_{\substack{g\in G\\ \dogo g\leq \rho}} b^{-s \dogo g} \geq \sum_{\substack{g\in G\\ \dogo g\leq \rho}} b^{-s\rho} = \infty .
\]
Since $s$ was arbitrary, we have $\delta_G = \infty$.
\end{proof}

\begin{remark}
Although the converse to Observation \ref{observationnotstronglydiscrete} holds in the Standard Case, it fails for infinite-dimensional algebraic hyperbolic spaces; see Example \ref{exampleinfinitepoincareBIM}.
\end{remark}

\begin{notation}
\label{notationpoincareset}
The Poincar\'e exponent and type can be conveniently combined into a single mathematical object, the \emph{Poincar\'e set}
\[
\Delta_G := \{s\geq 0 : \Sigma_s(G) = \infty\} = \begin{cases}
[0,\delta_G] & \text{$G$ is of divergence type}\\
\CO 0{\delta_G} & \text{$G$ is of convergence type}\\
\Rplus & \delta_G = \infty
\end{cases}.
\]
\end{notation}

\section{The modified Poincar\'e exponent of a semigroup}
\label{subsectionmodified}

From a certain perspective, Observation \ref{observationnotstronglydiscrete} indicates a flaw in the Poincar\'e exponent: If $G\leq\Isom(X)$ is not strongly discrete, then the Poincar\'e exponent is always infinity even though there may be more geometric information to capture. In this section we introduce a modification of the Poincar\'e exponent which agrees with the Poincar\'e exponent in the case where $G$ is strongly discrete, but can be finite even if $G$ is not strongly discrete.

We begin by defining the modified Poincar\'e exponent of a locally compact group $G\leq\Isom(X)$. Let $\mu$ be a Haar measure on $G$, and for each $s$ consider the \emph{Poincar\'e integral}
\begin{equation}
\label{IsG}
I_s(G) := \int b^{-s\dogo g}\;\dee\mu(g).
\end{equation}
\begin{definition}
\label{definitionmodified1}
The \emph{modified Poincar\'e exponent} of a locally compact group $G\leq\Isom(X)$ is the number
\[
\w\delta_G = \w\delta(G) := \inf\{s\geq 0: I_s(G) < \infty\},
\]
where $I_s(G)$ is defined by \eqref{IsG}.
\end{definition}
\begin{example}
\label{exampleLie}
Let $X = \H^d$ for some $2\leq d < \infty$, and let $G\leq\Isom(X)$ be a positive-dimensional Lie subgroup. Then $G$ is locally compact, but not strongly discrete. Although the Poincar\'e series diverges for every $s$, the exponent of convergence of the Poincar\'e integral (or ``modified Poincar\'e exponent'') is equal to the Hausdorff dimension of the limit set of $G$ (Theorem \ref{theorembishopjonesmodified} below), and so in particular the Poincar\'e integral converges whenever $s > d - 1$.
\end{example}

We now proceed to generalize Definition \ref{definitionmodified1} to the case where $G\leq\Isom(X)$ is not necessarily locally compact. Fix $\rho > 0$, and consider a maximal $\rho$-separated\Footnote{Here, as usual, a \emph{$\rho$-separated} subset of a metric space $X$ is a set $S\subset X$ such that $\dist(x,y)\geq \rho$ for any distinct $x,y\in S$. The existence of a maximal $\rho$-separated subset of any metric space is guaranteed by Zorn's lemma.\label{footnoterhoseparated}} subset $S_\rho\subset G(\zero)$. Then we have
\[
\bigcup_{x\in S_\rho} B(x,\rho/2) \subset G(\zero) \subset \bigcup_{x\in S_\rho} B(x,\rho),
\]
and the former union is disjoint. Now suppose that $G$ is in fact locally compact, and let $\nu$ denote the image of Haar measure on $G$ under the map $g\mapsto g(\zero)$. Then if $f$ is a positive function on $X$ whose logarithm is uniformly continuous, we have
\[
\sum_{x\in S_\rho} f(x)
\asymp_{\times,\rho,f} \sum_{x\in S_\rho} \int_{B(x,\rho/2)} f\;\dee\nu
\leq \int f\;\dee\nu
\leq \sum_{x\in S_\rho} \int_{B(x,\rho)} f\;\dee\nu
\asymp_{\times,\rho,f} \sum_{x\in S_\rho} f(x).
\]
Thus in some sense, the counting measure on $S_\rho$ is a good approximation to the measure $\nu$. In particular, taking $f(x) = b^{-\dox x}$ gives
\[
I_s(G) \asymp_{\times,\rho} \sum_{x\in S_\rho} b^{-\dox x}.
\]
Thus the integral $I_s(G)$ converges if and only if the series $\sum_{x\in S_\rho} b^{-\dox x}$ converges. But the latter series is well-defined even if $G$ is not locally compact. This discussion shows that the definition of the ``modified Poincar\'e exponent'' given in Definition \ref{definitionmodified1} agrees with the following definition:

\begin{definition}
\label{definitionmodifiedexponent}
Fix $G\prec\Isom(X)$.
\begin{itemize}
\item For each set $S\subset X$ and $s\geq 0$, let
\begin{align*}
\Sigma_s(S) &= \sum_{x\in S} b^{-s\dox x}\\
\Delta(S) &= \{s\geq 0: \Sigma_s(S) = \infty\}\\
\delta(S) &= \sup\Delta(S).
\end{align*}
\item Let
\begin{equation}
\label{modifiedpoincaredef}
\w\Delta_G = \bigcap_{\rho > 0} \bigcap_{S_\rho} \Delta(S_\rho),
\end{equation}
where the second intersection is taken over all maximal $\rho$-separated sets $S_\rho$.
\item The number $\w\delta_G = \sup\w\Delta_G$ is called the \emph{modified Poincar\'e exponent} of $G$. If $\w\delta_G\in\w\Delta_G$, we say that $G$ is of \emph{generalized divergence type},\Footnote{We use the adjective ``generalized'' rather than ``modified'' because all groups of convergence/divergence type are also of generalized convergence/divergence type; see Corollary \ref{corollarygeneralizedtypes} below.} while if $\w\delta_G\in\Rplus\butnot\Delta_G$, we say that $G$ is of \emph{generalized convergence type}. Note that if $\w\delta_G = \infty$, then $G$ is neither of generalized convergence type nor of generalized divergence type.
\end{itemize}
\end{definition}

The basic properties of the modified Poincar\'e exponent are summarized as follows:

\begin{proposition}
\label{propositionbasicmodified}
Fix $G\prec\Isom(X)$.
\begin{itemize}
\item[(i)] $\w\Delta_G \subset \Delta_G$. (In particular $\w\delta_G\leq\delta_G$.)
\item[(ii)] If $G$ satisfies
\begin{equation}
\label{SD3}
\sup_{x\in X} \#\{g\in G:\dist(g(\zero),x)\leq \rho\} < \infty \all \rho > 0,
\end{equation}
then $\w\Delta_G = \Delta_G$. (In particular $\w\delta_G = \delta_G$.)
\item[(iii)] If $\w\delta_G < \infty$, then there exist $\rho > 0$ and a maximal $\rho$-separated set $S_\rho\subset G(\zero)$ such that $\#(S_\rho\cap B) < \infty$ for every bounded set $B$.
\item[(iv)] For all $\rho > 0$ sufficiently large and for every maximal $\rho$-separated set $S_\rho\subset G(\zero)$, we have $\Delta(S_\rho) = \w\Delta_G$. (In particular $\delta(S_\rho) = \w\delta(G)$.)
\end{itemize}
\end{proposition}
\begin{remark}
\label{remarkstronglydiscrete}
If $G$ is a group, then it is clear that \eqref{SD3} is equivalent to the assertion that $G$ is strongly discrete. If $G$ is not a group, then by analogy we will say that $G$ is \emph{strongly discrete} if \eqref{SD3} holds. (Recall that in Chapter \ref{sectiondiscreteness}, the various notions of discreteness are defined only for groups.)
\end{remark}

\begin{proof}[Proof of Proposition \ref{propositionbasicmodified}]
~
\begin{itemize}
\item[(i)] Indeed, for every $s\geq 0$, $\rho > 0$, and maximal $\rho$-separated set $S_\rho$ we have
\[
\Sigma_s(S_\rho) \leq \Sigma_s(G) \text{ and thus }\w\Delta(G) \subset \Delta(S_\rho) \subset \Delta(G).
\]
\item[(ii)] Fix $\rho > 0$, and let $S_\rho\subset G(\zero)$ be a maximal $\rho$-separated set. For every $x\in G(\zero)$ there exists $y_x\in S_\rho$ with $\dist(x,y_x) \leq \rho$. Then for each $y\in S_\rho$, we have
\[
\#\{ x\in G(\zero) : y_x = y \} \leq M_\rho,
\]
where $M_\rho$ is the value of the supremum \eqref{SD3}. Therefore for each $s\geq 0$ we have
\[
\Sigma_s(G)
= \sum_{x\in G(\zero)} b^{-s \dox x}
\asymp_\times \sum_{x\in G(\zero)} b^{-s\dox{y_x}}
\leq M_\rho \sum_{y\in S_\rho} b^{-s \dox y}
= M_\rho\Sigma_s(S_\rho).
\]
In particular, $\Sigma_s(G) < \infty$ if and only if $\Sigma_s(S_\rho) < \infty$, i.e. $\Delta(G) = \Delta(S_\rho)$. Intersecting over $\rho > 0$ and $S_\rho\subset G(\zero)$ yields $\Delta(G) = \w\Delta(G)$.
\item[(iii)] Take $\rho$ and $S_\rho$ such that $\delta(S_\rho) < \infty$.
\end{itemize}

Before proving (iv), we need a lemma:
\begin{lemma}
\label{lemma2tau}
Fix $\rho_1,\rho_2 > 0$ with $\rho_2\geq 2\rho_1$. Let $S_1\subset G(\zero)$ be a $\rho_1$-net,\Footnote{Here, as usual, a \emph{$\rho$-net} in a metric space $X$ is a subset $S\subset X$ such that $X = \thicken_\rho(S)$. Note that every maximal $\rho$-separated set is a $\rho$-net (but not conversely). \label{footnoterhonet}} and let $S_2\subset G(\zero)$ be a $\rho_2$-separated set. Then
\begin{equation}
\label{2tau}
\Delta(S_2) \subset \Delta(S_1).
\end{equation}
\end{lemma}
\begin{subproof}
Since $S_1$ is a $\rho_1$-net, for every $y\in S_2$, there exists $x_y\in S_1$ with $\dist(y,x_y) < \rho_1$. If $x_y = x_z$ for some $y,z\in S_2$, then $\dist(y,z) < 2\rho_1 \leq \rho_2$ and since $S_2$ is $\rho_2$-separated we have $y = z$. Thus the map $y\mapsto x_y$ is injective. It follows that for every $s\geq 0$, we have
\[
\Sigma_s(S_2)
= \sum_{y\in S_2} b^{-s \dox y}
\asymp_\times \sum_{y\in S_2} b^{-s \dox{x_y}}
\leq \sum_{x\in S_1} b^{-s \dox x}
= \Sigma_s(S_1),
\]
demonstrating \eqref{2tau}.
\end{subproof}

\begin{itemize}
\item[(iv)] The statement is trivial if $\w\delta_G = \infty$. So suppose that $\w\delta_G < \infty$, and let $\rho$, $S_\rho$ be as in (iii). Fix $\rho'\geq 2\rho$ and a maximal $\rho'$-separated set $S_{\rho'}\subset G(\zero)$, and we will show that $\Delta(S_{\rho'}) = \w\Delta_G$. The inclusion $\supset$ follows by definition. To prove the reverse direction, fix $\rho'' > 0$ and a maximal $\rho''$-separated set $S_{\rho''}$, and we will show that $\Delta(S_{\rho'})\subset \Delta(S_{\rho''})$.

Let $F = S_\rho\cap B(\zero,\rho'' + \rho)$; then $\#(F) < \infty$. We then set 
\[
S_\rho' := \bigcup_{x\in S_{\rho''}} g_x(F)
\] 
where for each $x\in S_{\rho''}$, $x = g_x(\zero)$. Then for all $s\geq 0$,
\begin{align*}
\Sigma_s(S_{\rho''}) = \sum_{x\in S_{\rho''}} b^{-s\dox x} 
&\asymp_\times \sum_{x\in S_{\rho''}} \sum_{y\in F} b^{-s\dox x}\\ 
&\asymp_\times \sum_{x\in S_{\rho''}} \sum_{y\in F} b^{-s\dox{g_x(y)}}\\ 
&= \Sigma_s(S_\rho')
\end{align*}
and therefore $\Delta(S_{\rho''}) = \Delta(S_\rho')$. But $S_\rho'$ is a $\rho$-net, so by Lemma \ref{lemma2tau}, we have $\Delta(S_{\rho'}) \subset \Delta(S_\rho')$. This completes the proof.
\end{itemize}
\end{proof}

Combining with Observation \ref{observationnotstronglydiscrete} yields the following:
\begin{corollary}
\label{corollarydeltawdelta}
Suppose that $G$ is a group. If $\Delta\neq\w\Delta$ then
\[
\w\delta < \delta = \infty.
\]
\end{corollary}

\begin{corollary}
\label{corollarygeneralizedtypes}
If a group $G$ is of convergence or divergence type, then it is also of generalized convergence or divergence type, respectively.
\end{corollary}

We will call a group $G\leq\Isom(X)$ \emph{Poincar\'e regular} if $\w\Delta_G = \Delta_G$, and \emph{Poincar\'e irregular} otherwise. A list of sufficient conditions for Poincar\'e regularity is given in Proposition \ref{propositionpoincareregular} below. Conversely, several examples of Poincar\'e irregular groups may be found in Section \ref{subsectionpoincareirregular}.


